\originalchapter{研究生作业十二：散射与软辐射}
题面据8.311作业12全部5题. 官方解答末页标有“Correct answers”, 仍有一处角谱幂次错误; 本章解答为AI独立推导并注明模型及勘误. 正频单边谱采用第10章傅里叶约定.
\originalsection{束缚电子与小介质球散射}
\begin{exercise}
原题1. (a) 电子被弹簧常数$K$束缚, 受圆偏振$E_0$驱动, 求散射截面; (b) 半径$a\ll\lambda$、相对介电率$\varepsilon_r$球受线偏振波, 求截面.
\end{exercise}
\begin{solution}
(a) $\omega_0^2=K/m$, 忽略阻尼, $m(\ddot{\rr}+\omega_0^2\rr)=-e\ee$, 圆周受迫振幅$a_r=eE_0/[m|\omega_0^2-\omega^2|]$, 恒模加速度$\omega^2a_r$. 圆偏振$\langle S\rangle=\eps cE_0^2$, 所以$\sigma=\sigma_T\omega^4/(\omega_0^2-\omega^2)^2$, $\sigma_T=e^4/(6\pi\eps^2m^2c^4)$. $\omega\gg\omega_0$恢复Thomson, 低频按$\omega^4$; 无阻尼共振没有有限周期稳态, 需加耗散方能求线宽.

(b) 第5章球边界联立给$\alpha=4\pi\eps a^3(\varepsilon_r-1)/(\varepsilon_r+2)$, $\pp=\alpha\ee$. 线偏振峰值$E_0$的$\langle\ddot p^2\rangle=\omega^4|\alpha|^2E_0^2/2$, 因而$P=\omega^4|\alpha|^2E_0^2/(12\pi\eps c^3)$, 除$\eps cE_0^2/2$得$\sigma=(8\pi/3)k^4a^6| (\varepsilon_r-1)/(\varepsilon_r+2)|^2$. 损耗材料此式为散射部分, 总消光还含吸收, 也要检查准静态极化率在共振附近的辐射修正.
\end{solution}
\Needspace{9\baselineskip}
\originalsection{短碰撞和金属屏蔽的低频谱}
\begin{exercise}
原题2. (a) 非相对论电荷短时间$\tau$由$\vct v_1$变$\vct v_2$, 求$\omega\tau\ll1$角谱; (b) 电子正入射金属, 进入后被完全屏蔽, 用镜像求真空侧谱.
\end{exercise}
\begin{solution}
(a) 第10章Parseval推导给$\dd^2W/(\dd\omega\dd\Omega)=q^2|\nn\times(\vct v_2-\vct v_1)|^2/(16\pi^3\eps c^3)$; 总谱$\dd W/\dd\omega=q^2|\Delta\vct v|^2/(6\pi^2\eps c^3)$, 用$\int\sin^2\vartheta\dd\Omega=8\pi/3$. 这是低频极限, 并非无限频率平坦谱.

(b) 表面法向入射, 实荷$q$与像荷$-q$速度相反, 横向电流振幅同号. 进入屏蔽后外部源消失, 有效$|\Delta\dot{\pp}|=2|q|v$, 所以真空上半空间$\dd^2W/(\dd\omega\dd\Omega)=q^2v^2\sin^2\vartheta/(4\pi^3\eps c^3)$. 只积分物理半空间给$\dd W/\dd\omega=q^2v^2/(3\pi^2\eps c^3)$, 是自由电荷停止全空间谱的两倍. 不能把像电荷下半空间也计作真实辐射.
\end{solution}
\originalsection{对称裂变的相干角谱}
\begin{exercise}
原题3. 静止核质量$M_0$、总物理电荷$Q$（原卷写$Z$）, 瞬间裂为两半, 各质量$M_1<M_0/2$、电荷$Q/2$, 以$\pm\vct v$分离. 求角谱和积分谱. 若$Z$解释原子序数则$Q=Ze$, 不默认原卷的$Z$一定是无量纲整数.
\end{exercise}
\begin{solution}
忽略辐射对末速度的一阶反馈, 质量能量给$\gamma=M_0/(2M_1)$, $\beta=\sqrt{1-(2M_1/M_0)^2}$. 取分离轴$z$, $x=\cos\vartheta$. 静止初源的软因子零, 两末源须先相加: $\sum (Q/2)\nn\times(\pm\vct v)/(1\mp\beta x)=Q\beta^2cx\,\nn\times\hat{\vct z}/(1-\beta^2x^2)$. 因而
\[
\frac{\dd^2W}{\dd\omega\dd\Omega}=\frac{Q^2}{16\pi^3\eps c}\frac{\beta^4x^2(1-x^2)}{(1-\beta^2x^2)^2}.
\]
官方修正版3(a)多写一个$\sin^2\vartheta$, 与横向因子及其同页角积分不一致, 此处独算为$\sin^2\vartheta\cos^2\vartheta$. 中央赤道与轴向皆零, 是对称源干涉的结果, 不能直接把两个单粒子功率相加.

积分$\dd\Omega=2\pi\dd x$. 令$A=\operatorname{artanh}\beta/\beta$, 用
$\int_{-1}^1x^2/(1-\beta^2x^2)^2\dd x=\beta^{-2}[(1-\beta^2)^{-1}-A]$以及
$\int_{-1}^1x^4/(1-\beta^2x^2)^2\dd x=\beta^{-4}[(1-\beta^2)^{-1}-3A+2]$, 相减得$J=(3-\beta^2)A-3$. 所以$\dd W/\dd\omega=Q^2J/(8\pi^2\eps c)$. 小$\beta$展开$J=4\beta^4/15+8\beta^6/35+\cdots$, 首阶电偶极抵消. 原理想瞬间模型无反应时标, 实际只应把该式作为低频软谱; 有限裂变时间与核结构限制高频, 不能积分常谱到无穷而宣称有限总能.
\end{solution}
\originalsection{弱横向偏转的辐射四动量}
\begin{exercise}
原题4. 相对论电子初速度$\vct v_0$垂直电容器电场$\ee$, 路程短、场弱且速度变化远小于$c$. (a) 初始共动惯性系用Larmor求辐射及与电子能动量变化的关系; (b) 变回实验室.
\end{exercise}
\begin{solution}
取板内实验室驻留时间$T$, 短路径还须$|q|ET/(mc)\ll1$使初始共动系中末运动非相对论, 才可作统一Larmor近似（仅实验室速度差小于$c$不充分）, $E'=\gamma_0E$, $T'=T/\gamma_0$, $a'=q\gamma_0E/m$. $W'_{\rm rad}=q^2a'^2T'/(6\pi\eps c^3)=q^4\gamma_0E^2T/(6\pi\eps m^2c^3)$. 共动低速电偶极角分布前后对称, 净$\pp'_{\rm rad}=0$至该阶. 电子变化仍须$\delta E'_e=W'_{\rm ext}-W'_{\rm rad}$、$\delta\pp'_e=\vct I'_{\rm ext}-\pp'_{\rm rad}$; 不是漏掉外功后直接$\delta E'_e=-W'_{\rm rad}$. 主横向冲量$\delta p'_\perp=qE'T'$, 无辐射主项动能$\delta p_\perp'^2/(2m)$.

辐射四动量反boost给$W_{\rm rad}=\gamma_0W'_{\rm rad}=q^4\gamma_0^2E^2T/(6\pi\eps m^2c^3)$, $\pp_{\rm rad}=\vct v_0W_{\rm rad}/c^2$. 所以平均功率$q^4\gamma_0^2E^2/(6\pi\eps m^2c^3)$, 亦可用横加速度$a=qE/(\gamma_0m)$代Liénard的$\gamma_0^4a^2$核对. 这与官方末页修正版一致, 仅在小偏转、慢能量损失范围内成立.
\end{solution}
\originalsection{相对论回旋与旋转电场模型}
\begin{exercise}
原题5. (a) 相对论电子垂直恒磁场回旋, 求瞬时辐射功率; (b) 圆偏振电场幅$E_0$驱动, 求圆偏振驱动辐射及截面, 比较非相对论与强度依赖, 并说明周期支和初静止支的不同.
\end{exercise}
\begin{solution}
(a) 横向圆运动$\vct v\cdot\vct a=0$, Liénard给$P=e^2\gamma^4a^2/(6\pi\eps c^3)$, $a=eBv/(\gamma m)$, 故$P=e^4B^2\gamma^2v^2/(6\pi\eps m^2c^3)$. 磁力不做功, 辐射使$\gamma$缓慢下降, 采用瞬时轨道近似, 不假定永久恒能圆轨道.

(b) 先给官方采用的无纵漂周期支, 再检查真实平面波的初值. 令$a_0=eE_0/(m\omega c)$, 圆周特解机械动量幅$p_a=eE_0/\omega$, $\gamma=\sqrt{1+a_0^2}$、$v=p_a/(\gamma m)$、$a=\omega v$. Liénard除圆偏振输入$\eps cE_0^2$得$\sigma=\sigma_T(1+a_0^2)$. 小参数恢复Thomson, 这个周期支随强度增强.

为说明它的实际适用域, 保留真实真空平面波磁场$\bb=\hat{\vct z}\times\ee/c$, 相位$\varphi=\omega(t-z/c)$. 取横势$\ee=-\partial_t\vct A_\perp$, 定义无量纲四速度$u^0=\gamma,u^i=\gamma v_i/c$. Lorentz式给$s=u^0-u^z>0$守恒, 横向正则量$\vct C=\vct u_\perp+q\vct A_\perp/(mc)$守恒. 归一化$(u^0)^2-u_z^2-u_\perp^2=1$给
\[
\vct u_\perp=\vct C-\frac{q\vct A_\perp(\varphi)}{mc},\quad
u^0=\frac{1+s^2+u_\perp^2}{2s},\quad
u^z=\frac{1-s^2+u_\perp^2}{2s},\quad
\frac{\dd\varphi}{\dd\tau}=\omega s.
\]
这就是忽略辐射反作用时含初态的完整平面波解: 初四速度决定$s,\vct C$, 位置由$\dd\rr/\dd\tau=c\vct u$、$\dd t/\dd\tau=u^0$积分; 各分量是已给相位正弦、余弦及其平方, 无未知动力学步骤.

圆偏振$|\partial_\varphi\vct A|=E_0/\omega$, 由于$\dot u^0=\dot u^z$, 四加速度范数中的两项相消, $-a_\mu a^\mu=c^2|\dot{\vct u}_\perp|^2=c^2\omega^2s^2a_0^2$. 所以总瞬时辐射$P=e^4E_0^2s^2/(6\pi\eps m^2c^3)$, 按本题实验室总功率除输入通量定义的$\sigma=\sigma_Ts^2$.

官方支可由$\vct C=0,s=\sqrt{1+a_0^2}$取得: $u_z=0,\gamma=s$, 速度始终垂直电场, $\vct v\times\bb=\hat{\vct z}(\vct v_\perp\cdot\ee)/c=0$, 因而它也可以是完整平面波的特殊周期解, 无需错误地说强场磁力永远无法为零. 但它有初始非零横向动量. 若电子从波前初静止进入, 则$s=1$, 存在纵漂; 圆偏振零横正则量支$u_z=a_0^2/2,\gamma=1+a_0^2/2$, 此功率/通量定义下$\sigma=\sigma_T$. 初值不同, 强度依赖也不同, 所以官方增强式不自动适用于任意初静止强波电子. 辐射反作用与有限包能量耗尽会改变这些外场轨道, 本题未给维持稳态的机制.
\end{solution}
